\section*{第 \ref{ch:b2:reaction-free-energy} 章——熵与反应自由能}

\begin{solution}{exo:b2:reaction-free-energy:1}
(a) 三摩尔气体变为液体：$\Delta_r S^\circ < 0$。(b) 生成了一种
气体：$> 0$。(c) 一个气体分子变为两个：$> 0$。(d) 溶液中的离子
形成晶体：$< 0$。(e) 熔化：$> 0$。(f) $\Delta\nu_{\text{气体}} = 0$：
数值较小，此处为正（由表中数据得 \qty{+20}{J/(K.mol)}，\ce{HCl}
其对称性低于 \ce{H2} 与 \ce{Cl2}）。
% ledger: s0:HCl_g, s0:H2_g, s0:Cl2_g
\end{solution}

\begin{solution}{exo:b2:reaction-free-energy:2}
$\Delta_r S^\circ = 2(192.77) - 191.61 - 3(130.68) = \qty{-198.1}{J/(K.mol)}$；
$\Delta_r G^\circ = -91.9 - 298.15 \times (-0.1981) = \qty{-32.8}{kJ/mol}$：
该反应在 \qty{298}{K} 下为\omterm{def:b2:reaction-free-energy:exergonic}{放能的}，尽管熵项（\qty{+59.1}{kJ/mol}）
起阻碍作用。
\end{solution}

\begin{solution}{exo:b2:reaction-free-energy:3}
$\Delta_r S^\circ = 69.95 - 130.68 - \tfrac12(205.15) = \qty{-163.3}{J/(K.mol)}$；
$\Delta_r G^\circ = -285.83 + 298.15 \times 0.1633 = \qty{-237.14}{kJ/mol}$。
该反应是液态水的生成反应（由处于参考态的单质生成
一摩尔液态水），因此其\omterm{def:b2:reaction-free-energy:reaction-gibbs}{标准反应吉布斯自由能}按定义
等于 $\Delta_f G^\circ(\ce{H2O}, \text{l})$；
查表所得数值与之相同。
\end{solution}

\begin{solution}{exo:b2:reaction-free-energy:4}
在 \qty{298}{K} 时约为 \qty{42}{J/(K.mol)}，在
\qty{1000}{K}（液态）时约为 \qty{87}{J/(K.mol)}。跃变：在 \qty{692.7}{K} 时为 \qty{10.6}{J/(K.mol)}，
在 \qty{1180.2}{K} 时为 \qty{97.7}{J/(K.mol)}。由此 $\Delta_{\text{熔化}}H = 10.6 \times
692.7 = \qty{7.3}{kJ/mol}$，$\Delta_{\text{汽化}}H = 97.7 \times 1180.2 =
\qty{115}{kJ/mol}$。
\end{solution}

\begin{solution}{exo:b2:reaction-free-energy:5}
$\Delta_r H^\circ = \qty{44.00}{kJ/mol}$，$\Delta_r S^\circ = 188.84 - 69.95 =
\qty{118.9}{J/(K.mol)}$，$\Delta_r G^\circ = 44.00 - 298.15 \times 0.1189 =
\qty{+8.55}{kJ/mol}$。在 \qty{298}{K} 下该变化不可逆（液态
水与 \qty{1}{bar} 下的蒸气不处于平衡），因此 $\Delta S \ne
Q/T$：$\Delta_r H^\circ/T = \qty{147.6}{J/(K.mol)}$。正的
$\Delta_r G^\circ$ 为正表示在 \qty{298}{K} 下水不会蒸发为
\qty{1}{bar} 的纯蒸气气氛，蒸气反而凝结。两相
在 $\Delta_r G^\circ = 0$ 时处于平衡：$T = 44\,004/118.9 =
\qty{370}{K}$，与正常沸点相差不到 \qty{3}{K}，
（\qty{373}{K}）相差不到 \qty{3}{K}：\omterm{def:b2:reaction-free-energy:ellingham-approximation}{埃林厄姆近似}在 \qty{75}{K} 的范围内是可靠的。
\end{solution}

\begin{solution}{exo:b2:reaction-free-energy:6}
$T_i = 57\,100/175.7 = \qty{325}{K}$。低于此温度时 $\Delta_r G^\circ > 0$，且
无色的 \ce{N2O4} 更为有利；冷却试管使气体向
\ce{N2O4} 方向移动，\ce{NO2} 的棕色变淡（与组成的定量联系
见\cref{ch:b2:equilibrium-shifts}）。
\end{solution}

\begin{solution}{exo:b2:reaction-free-energy:7}
$\Delta_r S^\circ = -\dd\Delta_r G^\circ/\dd T = -(-4.0 + 12.0)/100 =
\qty{-0.080}{kJ/(K.mol)} = \qty{-80}{J/(K.mol)}$，且 $\Delta_r H^\circ =
\Delta_r G^\circ + T\Delta_r S^\circ = -12.0 + 300 \times (-0.080) =
\qty{-36.0}{kJ/mol}$。校验：$\Delta_r G^\circ/T$ 由 $-0.0400$ 变为
$\qty{-0.0100}{kJ/(K.mol)}$，斜率为 $3.00 \times 10^{-4}$ 每开尔文，等于
$-\Delta_r H^\circ/T^2 = 36.0/346^2 = 3.00 \times 10^{-4}$，此为
几何平均温度 $\sqrt{300 \times 400} = \qty{346}{K}$ 处的值。
\end{solution}

\begin{solution}{exo:b2:reaction-free-energy:8}
$\Delta_r S^\circ = 186.25 - 5.74 - 2(130.68) = \qty{-80.8}{J/(K.mol)}$，
$\Delta_r G^\circ = -74.87 + 298.15 \times 0.0808 = \qty{-50.8}{kJ/mol}$，
即表列甲烷的 $\Delta_f G^\circ$。两项符号相同，
甲烷在高于 $T_i = 74\,870/80.8 \approx
\qty{926}{K}$ 时变得不稳定（\omterm{def:b2:reaction-free-energy:ellingham-approximation}{埃林厄姆近似}）：热甲烷裂解为碳和
氢，这是同时获得两者的途径。
\end{solution}

\begin{solution}{exo:b2:reaction-free-energy:9}
取代：\ce{CH3CHBrCH3 + OH- -> CH3CH(OH)CH3 + Br-}；消除：
\ce{CH3CHBrCH3 + OH- -> CH2=CHCH3 + H2O + Br-}。两者的\omterm{def:b2:reaction-free-energy:reaction-gibbs}{标准反应吉布斯自由能}之差
为 $\Delta(\Delta_r G^\circ) = 40 - T \times
0.090$，在高于 $40\,000/90 \approx \qty{440}{K}$ 时为负。消除反应
比取代多生成一个分子，其熵更大，熵项
随 $T$ 成正比增长。
\end{solution}

\begin{solution}{exo:b2:reaction-free-energy:10}
$N_A$ 个分子中每一个都有 2 种取向：$W = 2^{N_A}$，故 $S =
k_B\ln 2^{N_A} = R\ln 2 = \qty{5.76}{J/(K.mol)}$。晶体并非
完全有序：在低温下，分子被冻结在其
随机取向之中，无法达到热力学第三定律所假设的
唯一有序排列。
\end{solution}

\begin{solution}{exo:b2:reaction-free-energy:11}
对 $\dd\Delta_r S^\circ/\dd T = c/T$ 积分：$\Delta_r S^\circ(T) = \Delta_r
S^\circ(T_0) + c\ln(T/T_0)$；由\omterm{thm:b2:reaction-enthalpy:kirchhoff}{基尔霍夫定律}可得 $\Delta_r H^\circ(T) = \Delta_r
H^\circ(T_0) + c(T - T_0)$；减去 $T\Delta_r S^\circ(T)$。对于石灰石，在
\qty{1100}{K} 时，\omterm{def:b2:reaction-free-energy:ellingham-approximation}{埃林厄姆近似}给出 $178.32 - 1100 \times 0.16064 = \qty{1.62}{kJ/mol}$；
修正项 $c[T - T_0 - T\ln(T/T_0)] = -0.00195 \times (801.85 - 1436.3) =
\qty{+1.24}{kJ/mol}$，故 $\Delta_r G^\circ = \qty{2.86}{kJ/mol}$。误差
\qty{1.2}{kJ/mol}，小于 $\Delta_r H^\circ$ 的百分之一，但在
\omterm{def:b2:reaction-free-energy:inversion-temperature}{转变温度}附近，它决定了符号。
\end{solution}

\begin{solution}{exo:b2:reaction-free-energy:12}
(a) $\Delta_r H^\circ = -763.2 + 944.0 = \qty{+180.8}{kJ/mol}$，$\Delta_r
S^\circ = 353.2 + 205.15 - 50.62 - 2(223.08) = \qty{+61.6}{J/(K.mol)}$；
$\Delta_r G^\circ(1000) = 180.8 - 61.6 = \qty{+119.2}{kJ/mol}$：\omterm{def:b2:reaction-free-energy:exergonic}{吸能的}，
不能发生氯化。(b) 对于 \ce{C + O2 -> CO2}，$\Delta_r G^\circ(1000) =
-393.5 - 2.9 = \qty{-396.4}{kJ/mol}$。两式相加得 \ce{TiO2(s) + 2Cl2(g) + C(s)
-> TiCl4(g) + CO2(g)}，其 $\Delta_r G^\circ = \qty{-277.2}{kJ/mol}$：碳
消耗了第一个反应生成的氧气，因而耦合反应
为强\omterm{def:b2:reaction-free-energy:exergonic}{放能的}。
\end{solution}

\begin{solution}{pb:b2:reaction-free-energy:1}
\textbf{1.} \ce{CaCO3(s) -> CaO(s) + CO2(g)}.
\textbf{2.} $\Delta_r H^\circ = -635.09 - 393.51 + 1206.92 =
\qty{+178.3}{kJ/mol}$：\omterm{def:b2:reaction-enthalpy:exothermic}{吸热的}。
\textbf{3.} $\Delta_r S^\circ = 39.75 + 213.79 - 92.9 = \qty{+160.6}{J/(K.mol)}$,
为正，因为由固体生成了气体。
\textbf{4.} $\Delta_r G^\circ = 178.32 - 298.15 \times 0.16064 =
\qty{+130.4}{kJ/mol}$.
\textbf{5.} 是：分解需要 $\Delta_r G < 0$，而这里它
为强正值；空气中二氧化碳的分压远低于 \qty{1}{bar}，
虽使 $\Delta_r G$ 降低，但远小于 \qty{130}{kJ/mol}（见下一章）。
\textbf{6.} $\Delta_r C^\circ_p = 42.80 + 37.13 - 81.88 =
\qty{-1.95}{J/(K.mol)}$，很小：$\Delta_r H^\circ$ 和 $\Delta_r S^\circ$
几乎不随 $T$ 变化。
\textbf{7.} $\Delta_r G^\circ(T) = 178.32 - 0.16064\,T$ (\unit{kJ/mol}).
\textbf{8.} 在 \qty{1000}{K} 时为 \qty{+17.7}{kJ/mol}；在
\qty{1300}{K} 时为 \qty{-30.5}{kJ/mol}。
\textbf{9.} $T_i = 178\,320/160.64 = \qty{1110}{K}$.
\textbf{10.} $\Delta_r G^\circ/T = 178.32/T - 0.16064$，其导数为
$-178.32/T^2 = -\Delta_r H^\circ/T^2$.
\textbf{11.} 在 $T_i$ 处，$c[T_i - T_0 - T_i\ln(T_i/T_0)] = -0.00195 \times
(811.9 - 1459.4) = \qty{+1.26}{kJ/mol}$；$\Delta_r G^\circ$ 归零需再升温
$1260/160.6 \approx \qty{8}{K}$，即约 \qty{1118}{K}，
相差不足百分之一。
\textbf{12.} 在 \qty{1000}{K} 时，$\Delta_r G > 0$：分解不能
进行（石灰会吸收二氧化碳）。在 \qty{1300}{K} 时，$\Delta_r G <
0$：石灰石分解。
\textbf{13.} $\Delta_r G$ 变为正：石灰吸收二氧化碳，
重新变为碳酸盐。
\textbf{14.} $\delta S_{\text{c}} = -\Delta_r G\,\dd\xi/T = 30\,506/1300 =
\qty{23.5}{J/K}$，每摩尔。
\textbf{15.} 是：\cref{ch:b2:chemical-potential}将说明
$\Delta_r G = \Delta_r G^\circ + RT\ln(p_{\ce{CO2}}/p^\circ)$，其值降低
当二氧化碳的压强低于 \qty{1}{bar} 时，使得
分解在 $T_i$ 以下即开始。
\textbf{16.} $n = 10^6/56.08 = \qty{1.783e4}{mol}$.
\textbf{17.} $1.783 \times 10^4 \times 178.32 = \qty{3.18e6}{kJ}$，即
每吨 \qty{3.18}{GJ}。
\textbf{18.} $1.783 \times 10^4 \times 44.01 = \qty{785}{kg}$ 的 \ce{CO2}。
\textbf{19.} 需供给的热量为 $3.18 \times 10^6/0.50 = \qty{6.36e6}{kJ}$，因此
$6.36 \times 10^6/802.3 = \qty{7.93e3}{mol}$ 的甲烷，\qty{127}{kg}。
\textbf{20.} $7.93 \times 10^3 \times 44.01 = \qty{349}{kg}$ 的 \ce{CO2}。
\textbf{21.} $785 + 349 = \qty{1134}{kg}$，每吨石灰，其中 69\,\%
来自石灰石。
\textbf{22.} 石灰石所释放的二氧化碳由化学计量关系决定，
无论如何加热石灰窑；只有燃料的那部分可以减少。
\textbf{23.} 在 \qty{1}{bar} 的二氧化碳压强下，石灰石在高于
其\omterm{def:b2:reaction-free-energy:inversion-temperature}{转变温度} $\boldsymbol{T_i \approx \qty{1.11e3}{K}}$ 时分解
（约 \qty{840}{\celsius}）。
\end{solution}
